% TOP_PROBLEM: {"id": 30006747, "title": "Superlinear Circuit Lower Bounds for Explicit Boolean Functions", "category_id": 15, "category": {"id": 15, "name": "computer_science", "display_name": "Computer Science", "description": "Computational complexity, algorithms, and theoretical CS.", "slug": "computer-science", "order_index": 15, "created_at": "2026-07-31T15:26:25.671Z"}, "status": "open"}
% FIRST_POSTED: null
% !TeX program = lualatex
% Compile twice: lualatex 30006747.tex
% All references and prose are embedded; no BibTeX database or external inputs.
\documentclass[11pt,a4paper]{article}
\usepackage[margin=24mm,headheight=14pt]{geometry}
\usepackage{fontspec}
\setmainfont{Latin Modern Roman}
\setsansfont{Latin Modern Sans}
\setmonofont{Latin Modern Mono}
\usepackage{amsmath,amssymb,mathtools}
\usepackage{unicode-math}
\setmathfont{Latin Modern Math}
\usepackage{microtype}
\usepackage{enumitem,xurl,needspace}
\usepackage[dvipsnames]{xcolor}
\usepackage{hyperref}
\hypersetup{unicode=true,colorlinks=true,linkcolor=MidnightBlue,urlcolor=MidnightBlue,
 pdftitle={Research attempt: Problem 30006747},pdfauthor={Research notebook},bookmarksnumbered=true}
\usepackage{bookmark}
\usepackage{tocloft}
\setlength{\cftsecnumwidth}{3em}
\cftsetpnumwidth{2.5em}
\cftsetrmarg{4em}
\usepackage{titlesec}
\titleformat{\section}{\Large\bfseries\raggedright}{\thesection}{0.7em}{}
\usepackage{fancyhdr}
\pagestyle{fancy}\fancyhf{}
\fancyhead[L]{\small\sffamily Open Problems Math | Research notebook}
\fancyhead[R]{\small Problem 30006747 | 27 September 2026}
\fancyfoot[C]{\thepage}
\setlength{\parindent}{0pt}
\setlength{\parskip}{3pt plus 1pt}
\setlength{\emergencystretch}{3em}
\setcounter{tocdepth}{1}
\setcounter{secnumdepth}{1}
\setlist[itemize]{leftmargin=3em,itemsep=5pt,topsep=4pt,parsep=0pt}
\allowdisplaybreaks
\newcommand{\N}{\mathbb{N}}
\newcommand{\Z}{\mathbb{Z}}
\newcommand{\Q}{\mathbb{Q}}
\newcommand{\R}{\mathbb{R}}
\newcommand{\C}{\mathbb{C}}
\newcommand{\F}{\mathbb{F}}
\newcommand{\A}{\mathbb{A}}
\newcommand{\PP}{\mathbb P}
\newcommand{\E}{\mathbb{E}}
\newcommand{\eps}{\varepsilon}
\newcommand{\dd}{\,\mathrm d}
\newcommand{\sref}[2]{\hyperlink{source:#1:#2}{[S#2]}}
\newcommand{\eref}[2]{\hyperlink{extra:#1:#2}{[E#2]}}
% Turn the next switch off for a shorter reading copy. The quotations preserve
% every exactTarget field, including qualifications omitted from a short summary.
\newif\ifshowcatalogue
\showcataloguetrue
\newcommand{\cataloguescope}[1]{\ifshowcatalogue
 \paragraph{Exact scope in the supplied catalogue.}
 \begin{quote}\small #1\end{quote}\fi}
\usepackage{etoolbox}
\AtBeginEnvironment{itemize}{\small}
\numberwithin{equation}{section}
% Standalone export. Original research content preserved.
\begin{document}
\setcounter{section}{0}
% ================================================================
% problemId: 30006747
% Canonical problem ID: 30006747
% exactTarget SHA-256: 5b1b5f059b26149eff04937bbbfdb646c16b6d6c3e2ba1a33f52ed1142484325
\clearpage
\section*{\texorpdfstring{Superlinear Circuit Lower Bounds for Explicit Boolean Functions}{Superlinear Circuit Lower Bounds for Explicit Boolean Functions}}\label{30006747}
\subsection{Definitions and mathematical statement}
Let $\mathcal B_2$ be the set of all binary Boolean gates, with unrestricted circuit depth and fan-out. Denote by $C_{\mathcal B_2}(f_n)$ the minimum number of gates in a single-output circuit for $f_n$. Seek one uniformly polynomial-time computable family satisfying
\[
 \lim_{n\to\infty}\frac{C_{\mathcal B_2}(f_n)}{n}=\infty.
\]
Equivalently, every fixed linear bound is violated for all sufficiently large input lengths. A lower bound $cn-o(n)$ for one fixed $c$, no matter how large, is not superlinear in this sense.
\par\smallskip\textit{Formulation and concept references:} \sref{30006747}{1}.
\cataloguescope{There exists a Boolean function family f\_n:\{0,1\}\textasciicircum{}n->\{0,1\}, n>=1, evaluated by a single deterministic polynomial-time algorithm on (1\textasciicircum{}n,x), whose minimum gate count C\_B2(f\_n) over single-output fan-in-two circuits using arbitrary binary Boolean gates satisfies C\_B2(f\_n)=omega(n): for every c>0 there is N such that C\_B2(f\_n)>c*n for every n>=N. The circuits have unrestricted depth and fan-out; no fixed linear constant counts as superlinear.}
\subsection{Short English statement}
Can an explicit polynomial-time Boolean function be proved to require more than every constant multiple of its input length in general circuit gates?
% BEGIN RESEARCH ADDITION 139
\subsection{Research attempt}

\paragraph{Current frontier.}
Li--Yang prove a $3.1n-o(n)$ lower bound for explicit affine dispersers over
the full binary basis. Their discussion of particular gate-elimination
limitations is not a barrier against all lower-bound methods. A fixed
constant cannot give the divergent ratio in this target. \sref{30006747}{1}
Restricted-model advances must remain separate: Rao's September 2026
manuscript gives exponential bounds for \emph{monotone} matching circuits,
not circuits with arbitrary binary gates. \eref{30006747}{1}

\paragraph{Attack route.}
Counting finds hard small truth tables but is too slow at the final input
length. Padding restores polynomial-time uniformity while destroying the
lower-bound ratio. We combine a slowly growing hard table with many
\emph{disjoint} input blocks and an XOR output. The missing step is a
quantitative direct-sum estimate for unrestricted circuits, not the
availability of hard tables at this small scale.

\paragraph{Attempt: construct a uniform hard-block family.}
Allow free constants; this strengthens circuits and thus makes a lower bound
in that model sufficient for the target. For $k\ge8$, set
\[
 s_k=\left\lfloor\frac{2^k}{16k}\right\rfloor.
\]
Topologically ordered circuits with $s_k$ gates have at most
\begin{equation}\label{30006747:count}
 (k+s_k+2)\bigl[16(k+s_k+2)^2\bigr]^{s_k}
\end{equation}
descriptions. Padding with projection gates includes all smaller circuits.
Since $k+s_k+2\le2^k$, the base-two logarithm is at most
\[
 k+\frac{2^k}{16k}(4+2k)<2^k.
\]
Thus fewer than the $2^{2^k}$ truth tables are represented. Let $h_k$ be
the lexicographically first missing table. This definition is effective:
enumerate descriptions, evaluate their tables, mark them and select the
first unmarked table. The running time and space are $2^{O(2^k)}$, and
$C_{\mathcal B_2}(h_k)>s_k$. It is not advice hiding an unspecified hard
function.

For sufficiently large $n$, choose
\[
 k(n)=\lfloor\log_2\log_2n\rfloor,
 \qquad t(n)=\lfloor n/k(n)\rfloor,
\]
and on disjoint $k$-bit blocks define
\begin{equation}\label{30006747:family}
 f_n(x)=\bigoplus_{i=1}^{t(n)}h_{k(n)}(x^{(i)}).
\end{equation}
Ignore the fewer than $k$ leftover bits and define the finitely many small
lengths arbitrarily. Because $2^k\le\log_2n$, construction of $h_k$ takes
$n^{O(1)}$ time with one fixed exponent; all subsequent table lookups and
XOR operations are polynomial too. The logarithms defining $k$ are computed
by integer bit lengths. Thus this is a genuinely uniform polynomial-time
family meeting the explicitness requirement.

\paragraph{The precise amplification estimate that would suffice.}
For these specific tables, suppose there is an absolute $\gamma>0$ such that
\begin{equation}\label{30006747:DS}
 C_{\mathcal B_2}\left(\bigoplus_{i=1}^{t}h_k(x^{(i)})\right)
       \ge\gamma t C_{\mathcal B_2}(h_k)
\end{equation}
uniformly at the required $k,t$. Then
\[
 \frac{C_{\mathcal B_2}(f_n)}n
 \ge c\frac{2^{k(n)}}{k(n)^2}
 \ge c'\frac{\log n}{(\log\log n)^2}\longrightarrow\infty.
\]
Floors change only constants at sufficiently large lengths. Unused bits
cannot lower this bound: fixing them in any circuit gives a circuit of no
greater size for the active-block function. This conclusion holds at every
sufficiently large $n$, not only a sparse subsequence.

Even $\gamma_k$ tending to zero would suffice if
$\gamma_k2^k/k^2\to\infty$. This quantitative relaxation is a possible
attack target. Neither version of \eqref{30006747:DS} is proved here. Its
implication for the original problem is a warning about its depth, not a
reason to infer it from disjointness of variables.

\paragraph{What separation of blocks actually proves.}
Consider the restricted architecture where each block enters its own disjoint
subcircuit producing one bit $g_i(x^{(i)})$, and every later gate sees only
these bits. If the output is \eqref{30006747:family}, fix all blocks except $i$.
The final function is unary in $g_i$ but must equal $h_k$ or $1-h_k$.
Since $h_k$ is nonconstant for large $k$, this unary function is nonconstant;
hence $g_i=h_k$ or $1-h_k$. Complementing costs at most one binary gate, so
the $i$th subcircuit has size at least $C(h_k)-1$. Summing disjoint gate
sets proves
\[
 C\ge t(C(h_k)-1)
\]
in this architecture. This verifies amplification when early cross-block
mixing is prohibited, not in the catalogue's unrestricted model.

\paragraph{Stress test: why separate restrictions do not add.}
In an unrestricted circuit, fixing all but one block gives an individual
lower bound, but repeating this for each block gives a maximum, not their
sum. A gate can survive many different restrictions. For example, on $t$ two-bit blocks first compute the XOR of all first
bits and the XOR of all second bits, then XOR those two results. The last
gate remains a genuinely binary XOR under every restriction leaving only
one block. Charging it once per block counts the same gate $t$ times.
The companion script checks this overlap directly; it is not merely
counting gates that become removable unary identities after restriction.

Replacing disjoint blocks by repeated identical inputs is worse:
$h_k(x)\oplus h_k(x)=0$. Disjoint variables avoid this cancellation but do
not force disjoint internal computations. Bare padding
$\widetilde f_n(x)=h_{k(n)}(x_1,\ldots,x_k)$ has a truth-table circuit of
size $O(2^k)=O(\log n)$, so cannot be the desired family. A proof must
control shared computation, not tacitly replace circuits by the
block-separated architecture.

\paragraph{Outcome.}
\textbf{Outcome: Conditional reduction.}
A fully uniform polynomial-time candidate family has been constructed, and
a specified disjoint-block direct-sum estimate would give an explicit
divergent lower-bound ratio. Amplification is proved only for the restricted
architecture; unrestricted mixing is the remaining gap. The counting
mechanism is standard and no new general circuit lower bound is claimed.
A solution of this weak superlinear problem would not alone imply
$\mathsf P\ne\mathsf{NP}$: the constructed family is in $\mathsf P$ and
could still have larger polynomial-size circuits.

% END RESEARCH ADDITION 139
\subsection{Sources}
\begin{itemize}\raggedright
\item[\textnormal{[S1]}]\hypertarget{source:30006747:1}{} 3.1n - o(n) Circuit Lower Bounds for Explicit Functions, ECCC TR21-023,20February2021.
\url{https://eccc.weizmann.ac.il/report/2021/023/download/}.
{\small\textit{Catalogue formulation source.}}
% BEGIN RESEARCH REFERENCES 139
\item[\textnormal{[E1]}]\hypertarget{extra:30006747:1}{} A. Rao,
\emph{Monotone circuit lower bounds from spread matchings}, author's
manuscript, 3 September 2026, Theorem 1 and Corollary 2.
\url{https://homes.cs.washington.edu/~anuprao/pubs/matchinglowerbound.pdf}.

% END RESEARCH REFERENCES 139
\end{itemize}

\end{document}
